\section{动机示例}

考虑 Large Lemma Miners 开源基准中的 \texttt{swap2}~\cite{peled2026artifact}。设计中三个寄存器的初始条件为 $\mathrm{Init}:(a,b,c)=(1,5,7)$，转移关系 $T$ 由 $a'=b$、$b'=c$、$c'=a$ 给出。令 $K=\{1,5,7\}$，待证明的目标属性为 $p: a\in K$。所有可达状态中，$a,b,c$ 始终是 $K$ 中三个值的排列，因此 $p$ 成立；但单独的 $p$ 不是一步归纳的：例如不可达的源状态 $(1,0,0)$ 满足 $p$，经过一步转移后 $a'=0$，违反 $p$。

为说明候选选择，设有三条在可达状态上均成立的候选引理，并按示意顺序排列：$\ell_1:a\ne b$、$\ell_2:b\in K$、$\ell_3:c\in K$。考虑组合 $S^*=\{\ell_2,\ell_3\}$：它不是该顺序的前缀。三个前缀 $\{\ell_1\}$、$\{\ell_1,\ell_2\}$、$\{\ell_1,\ell_2,\ell_3\}$ 分别可由源状态 $(1,0,0)$、$(1,5,0)$、$(1,5,5)$ 构造一步反例：相应后继依次违反 $p$、$\ell_2$、$\ell_1$。因此，只检查这些前缀无法检验到 $S^*$。下面用一步归纳检查验证这个非前缀组合，并展示反例如何引导组合修复。

下列检查中，$I_0$ 和 $I^*$ 分别表示目标属性与相应候选组合的合取；对任意状态谓词 $I$，$I'$ 表示将其当前状态变量替换为后继状态变量得到的谓词。

\noindent\begin{minipage}{\columnwidth}
\hrule\smallskip
\noindent\textbf{示例检查：从 $S_0$ 到 $S^*$}\par
\smallskip\hrule\smallskip
{\small
\begin{enumerate}
  \setlength{\topsep}{2pt}
  \setlength{\itemsep}{1pt}
  \setlength{\parsep}{0pt}
  \item 取 $S_0=\{\ell_2\}$，令 $I_0=p\land\ell_2$。$\mathrm{Init}\land\neg I_0$ 不可满足，故初始性成立。
  \item 步骤查询 $I_0\land T\land\neg I_0'$ 可满足：$(1,5,0)\to(5,0,1)$。后继中 $b'=0\notin K$，故归纳保持失败。
  \item 由 $b'=c$ 可知，$\ell_2'$ 失败源于源状态的 $c=0$。加入 $\ell_3:c\in K$ 排除此状态，令 $I^*=p\land\ell_2\land\ell_3$。
  \item 将前述两项查询中的 $I_0$ 替换为 $I^*$，结果均不可满足；故 $I^*$ 是归纳不变式，且 $I^*\Rightarrow p$。
\end{enumerate}
}
\hrule
\end{minipage}\par\smallskip

这说明候选组合的归纳步骤失败未必应导致全局删除其中的引理；它也可能提示需要另一条候选提供联合支持。
